\begin{proof}[Proof of \cref{obj:lem:two-point-transfer}]
All expectations below are nonnegative extended integrals. We establish the lower bounds for the extended minimax values and then use their finiteness to pass to the real-valued criteria.

\begin{enumerate}
\item
By \cref{obj:lem:legal-cancellation}, the alternatives satisfy
\[
P^+_{b,m},P^-_{b,m}\in\mathcal M_\beta(\sigma),
\qquad
\Delta_{\mathrm{wit}}
:=
\left|\theta(P^+_{b,m})-\theta(P^-_{b,m})\right|
=
2\kappa\left(\frac{b}{m^2}\right)^\beta.
\]
Include the independent procedure seed in the experiment. Define the sample and input spaces, the seed law, and the two sample and experiment laws by
\[
\begin{gathered}
\Omega_{\mathrm{sam},n}
=
\bigl(\mathbb R\times\{0,1\}\times\{0,1\}\bigr)^n,
\qquad
\Omega_n=\Omega_{\mathrm{sam},n}\times[0,1],
\qquad
\lambda_U=\operatorname{Uniform}[0,1],
\\
\mathsf S_\pm
=
\bigl((P^\pm_{b,m})_{\mathrm{obs}}\bigr)^{\otimes n},
\qquad
\mathsf P_\pm
=
\mathsf S_\pm\otimes\lambda_U
=
\mathbb P_{P^\pm_{b,m},n}.
\end{gathered}
\]
These are probability laws on their respective spaces.

For a measurable event \(\mathcal E_{\mathrm{seed}}\subseteq\Omega_n\), define its seed-section probability by
\[
\varphi_{\mathrm{seed}}(\mathbf s)
=
\lambda_U\{u\in[0,1]:(\mathbf s,u)\in\mathcal E_{\mathrm{seed}}\},
\qquad
\mathbf s\in\Omega_{\mathrm{sam},n}.
\]
This function is measurable and takes values in \([0,1]\). Product integration and the layer-cake identity give
\[
\begin{aligned}
\left|
\mathsf P_+(\mathcal E_{\mathrm{seed}})
-
\mathsf P_-(\mathcal E_{\mathrm{seed}})
\right|
&=
\left|
\int\varphi_{\mathrm{seed}}\,d\mathsf S_+
-
\int\varphi_{\mathrm{seed}}\,d\mathsf S_-
\right|
\\
&=
\left|
\int_0^1
\left(
\mathsf S_+\{\varphi_{\mathrm{seed}}\ge t\}
-
\mathsf S_-\{\varphi_{\mathrm{seed}}\ge t\}
\right)\,dt
\right|
\\
&\le
\int_0^1\operatorname{TV}(\mathsf S_+,\mathsf S_-)\,dt
=
\operatorname{TV}(\mathsf S_+,\mathsf S_-).
\end{aligned}
\]
Taking the supremum over measurable events therefore yields
\[
\operatorname{TV}(\mathsf P_+,\mathsf P_-)
\le
\operatorname{TV}(\mathsf S_+,\mathsf S_-)
\le\frac15.
\]
Finally, define the witness amplitude by
\[
\delta_{\mathrm{wit}}
:=
\kappa\left(\frac{b}{m^2}\right)^\beta
\ge0,
\qquad
\Delta_{\mathrm{wit}}=2\delta_{\mathrm{wit}}.
\]
Its nonnegativity follows from \(\kappa=1/100\) and \(b/m^2>0\).
% lean: CausalSmith.Stat.RdTruesideNoiseFrontier.experiment_tv_le

\item
For every measurable estimator \(T:\Omega_n\to\mathbb R\) and every \(P\in\mathcal M_\beta(\sigma)\), define
\[
\mathfrak r_{\mathrm{abs}}(T,P)
:=
\int_{\Omega_n}|T(z)-\theta(P)|\,d\mathbb P_{P,n}(z)
\in[0,\infty].
\]
The extended minimax absolute risk is
\[
R^{\mathrm{ext}}_\beta(n,\sigma)
:=
\inf_{\substack{T:\Omega_n\to\mathbb R\\T\ \mathrm{measurable}}}
\ \sup_{P\in\mathcal M_\beta(\sigma)}
\mathfrak r_{\mathrm{abs}}(T,P).
\]
For a law in the benchmark class, \cref{obj:ass:mean-bounds} gives
\[
\frac14\le\mu_d(P,0)\le\frac34
\quad(d\in\{0,1\}),
\qquad
-\frac12
\le
\theta(P)=\mu_1(P,0)-\mu_0(P,0)
\le\frac12.
\]
Hence the constant estimator
\[
T_{\mathrm{zero}}(z):=0
\qquad(z\in\Omega_n)
\]
satisfies
\[
\mathfrak r_{\mathrm{abs}}(T_{\mathrm{zero}},P)
=
|\theta(P)|
\le\frac12.
\]
Consequently,
\[
0\le R^{\mathrm{ext}}_\beta(n,\sigma)\le\frac12<\infty.
\]
Thus its conversion to a real number preserves its value and order.
% lean: CausalSmith.Stat.RdTruesideNoiseFrontier.risk_value_ne_top

\item
We next obtain the estimator-wise sum bound. Define the dominating measure, its two Radon--Nikodym densities, and their difference by
\[
\xi_{\mathrm{pair}}:=\mathsf P_++\mathsf P_-,
\qquad
p_\pm(z):=\frac{d\mathsf P_\pm}{d\xi_{\mathrm{pair}}}(z),
\qquad
\eta_{\mathrm{pair}}(z):=p_+(z)-p_-(z),
\quad z\in\Omega_n.
\]
Choose finite nonnegative measurable versions of the densities, assigning zero on any exceptional null set. Define the common measure by
\[
\mathsf C_{\mathrm{pair}}(\mathcal E)
:=
\int_{\mathcal E}\min\{p_+(z),p_-(z)\}\,d\xi_{\mathrm{pair}}(z)
\]
for every measurable \(\mathcal E\subseteq\Omega_n\). The pointwise bounds on the minimum imply
\[
\mathsf C_{\mathrm{pair}}\le\mathsf P_+,
\qquad
\mathsf C_{\mathrm{pair}}\le\mathsf P_-.
\]

Considering separately \(p_+\le p_-\) and \(p_-\le p_+\) gives
\[
\min\{p_+,p_-\}
=
\frac{p_++p_--|p_+-p_-|}{2}.
\]
Moreover,
\[
\int p_+\,d\xi_{\mathrm{pair}}
=
\int p_-\,d\xi_{\mathrm{pair}}
=1,
\qquad
\int\eta_{\mathrm{pair}}\,d\xi_{\mathrm{pair}}=0.
\]
To identify the total variation term, define
\[
\mathcal E_{\mathrm{pos}}
:=
\{z\in\Omega_n:\eta_{\mathrm{pair}}(z)\ge0\}.
\]
The zero integral of \(\eta_{\mathrm{pair}}\) and its signs on this event and its complement imply
\[
\begin{aligned}
\int_{\mathcal E_{\mathrm{pos}}}\eta_{\mathrm{pair}}\,d\xi_{\mathrm{pair}}
&=
-\int_{\mathcal E_{\mathrm{pos}}^c}\eta_{\mathrm{pair}}\,d\xi_{\mathrm{pair}}
\\
&=
\frac12\int|\eta_{\mathrm{pair}}|\,d\xi_{\mathrm{pair}}.
\end{aligned}
\]
For any measurable \(\mathcal E\subseteq\Omega_n\), splitting the zero total integral over the event and its complement gives
\[
\int_{\mathcal E}\eta_{\mathrm{pair}}\,d\xi_{\mathrm{pair}}
+
\int_{\mathcal E^c}\eta_{\mathrm{pair}}\,d\xi_{\mathrm{pair}}
=0,
\qquad
\int_{\mathcal E^c}\eta_{\mathrm{pair}}\,d\xi_{\mathrm{pair}}
=
-\int_{\mathcal E}\eta_{\mathrm{pair}}\,d\xi_{\mathrm{pair}}.
\]
The absolute-density integral splits as
\[
\int|\eta_{\mathrm{pair}}|\,d\xi_{\mathrm{pair}}
=
\int_{\mathcal E}|\eta_{\mathrm{pair}}|\,d\xi_{\mathrm{pair}}
+
\int_{\mathcal E^c}|\eta_{\mathrm{pair}}|\,d\xi_{\mathrm{pair}}.
\]
Integrating the pointwise bounds on the signed density and its negative over each piece yields
\[
\begin{aligned}
\int_{\mathcal E}\eta_{\mathrm{pair}}\,d\xi_{\mathrm{pair}}
&\le\int_{\mathcal E}|\eta_{\mathrm{pair}}|\,d\xi_{\mathrm{pair}},
\\
\int_{\mathcal E^c}\eta_{\mathrm{pair}}\,d\xi_{\mathrm{pair}}
&\le\int_{\mathcal E^c}|\eta_{\mathrm{pair}}|\,d\xi_{\mathrm{pair}},
\\
-\int_{\mathcal E}\eta_{\mathrm{pair}}\,d\xi_{\mathrm{pair}}
&\le\int_{\mathcal E}|\eta_{\mathrm{pair}}|\,d\xi_{\mathrm{pair}},
\\
-\int_{\mathcal E^c}\eta_{\mathrm{pair}}\,d\xi_{\mathrm{pair}}
&\le\int_{\mathcal E^c}|\eta_{\mathrm{pair}}|\,d\xi_{\mathrm{pair}}.
\end{aligned}
\]
Adding the third and second inequalities, and then the first and fourth, and using the complement identity gives
\[
\begin{aligned}
-2\int_{\mathcal E}\eta_{\mathrm{pair}}\,d\xi_{\mathrm{pair}}
&\le\int|\eta_{\mathrm{pair}}|\,d\xi_{\mathrm{pair}},
\\
2\int_{\mathcal E}\eta_{\mathrm{pair}}\,d\xi_{\mathrm{pair}}
&\le\int|\eta_{\mathrm{pair}}|\,d\xi_{\mathrm{pair}}.
\end{aligned}
\]
Therefore
\[
\left|\mathsf P_+(\mathcal E)-\mathsf P_-(\mathcal E)\right|
=
\left|\int_{\mathcal E}\eta_{\mathrm{pair}}\,d\xi_{\mathrm{pair}}\right|
\le
\frac12\int|\eta_{\mathrm{pair}}|\,d\xi_{\mathrm{pair}},
\]
with equality for \(\mathcal E=\mathcal E_{\mathrm{pos}}\). It follows that
\[
\operatorname{TV}(\mathsf P_+,\mathsf P_-)
=
\frac12\int|p_+-p_-|\,d\xi_{\mathrm{pair}},
\]
and hence
\[
\begin{aligned}
\mathsf C_{\mathrm{pair}}(\Omega_n)
&=
\int\min\{p_+,p_-\}\,d\xi_{\mathrm{pair}}
\\
&=
\frac{1+1-\int|p_+-p_-|\,d\xi_{\mathrm{pair}}}{2}
=
1-\operatorname{TV}(\mathsf P_+,\mathsf P_-).
\end{aligned}
\]

Now fix any measurable estimator \(T:\Omega_n\to\mathbb R\). The triangle inequality gives, pointwise,
\[
\Delta_{\mathrm{wit}}
\le
|T(z)-\theta(P^+_{b,m})|
+
|T(z)-\theta(P^-_{b,m})|.
\]
Integrating against the common measure and using its domination by each experiment law yields
\[
\begin{aligned}
\Delta_{\mathrm{wit}}
\bigl(1-\operatorname{TV}(\mathsf P_+,\mathsf P_-)\bigr)
&=
\int\Delta_{\mathrm{wit}}\,d\mathsf C_{\mathrm{pair}}
\\
&\le
\int
\left(
|T-\theta(P^+_{b,m})|
+
|T-\theta(P^-_{b,m})|
\right)\,d\mathsf C_{\mathrm{pair}}
\\
&=
\int|T-\theta(P^+_{b,m})|\,d\mathsf C_{\mathrm{pair}}
+
\int|T-\theta(P^-_{b,m})|\,d\mathsf C_{\mathrm{pair}}
\\
&\le
\mathfrak r_{\mathrm{abs}}(T,P^+_{b,m})
+
\mathfrak r_{\mathrm{abs}}(T,P^-_{b,m}).
\end{aligned}
\]
These are inequalities of nonnegative extended integrals, including when an individual risk is infinite.
% lean: CausalSmith.Stat.RdTruesideNoiseFrontier.twoPoint_absRisk_sum

\item
Define the candidate estimation lower bound by
\[
c_{\mathrm{est}}:=\frac45\,\delta_{\mathrm{wit}}\ge0.
\]
The experiment-law proximity and the preceding sum bound imply
\[
2c_{\mathrm{est}}
=
\frac45\,\Delta_{\mathrm{wit}}
\le
\Delta_{\mathrm{wit}}
\bigl(1-\operatorname{TV}(\mathsf P_+,\mathsf P_-)\bigr)
\le
\mathfrak r_{\mathrm{abs}}(T,P^+_{b,m})
+
\mathfrak r_{\mathrm{abs}}(T,P^-_{b,m}).
\]
If the maximum of these two risks were less than \(c_{\mathrm{est}}\), both risks would be finite and less than \(c_{\mathrm{est}}\), giving
\[
\mathfrak r_{\mathrm{abs}}(T,P^+_{b,m})
+
\mathfrak r_{\mathrm{abs}}(T,P^-_{b,m})
<
2c_{\mathrm{est}},
\]
a contradiction. Thus
\[
c_{\mathrm{est}}
\le
\max\left\{
\mathfrak r_{\mathrm{abs}}(T,P^+_{b,m}),
\mathfrak r_{\mathrm{abs}}(T,P^-_{b,m})
\right\}
\le
\sup_{P\in\mathcal M_\beta(\sigma)}
\mathfrak r_{\mathrm{abs}}(T,P).
\]
Taking the infimum over measurable estimators gives
\[
R^{\mathrm{ext}}_\beta(n,\sigma)\ge c_{\mathrm{est}}.
\]
The established finiteness permits conversion to the real-valued criterion, so
\[
R_\beta(n,\sigma)
\ge
\frac45\,\kappa\left(\frac{b}{m^2}\right)^\beta.
\]
% lean: Causalean.Stat.le_minimaxValueENNReal_of_two_point

\item
For an honest interval procedure \(I\in\mathcal I_{\beta,n,\sigma}\), with endpoints \(\ell,\upsilon\) as in \cref{obj:def:honest-class}, and a benchmark law \(P\), define
\[
\mathfrak r_{\mathrm{len}}(I,P)
:=
\int_{\Omega_n}
\max\{\upsilon(z)-\ell(z),0\}\,d\mathbb P_{P,n}(z)
\in[0,\infty].
\]
Define the extended minimax length by
\[
\mathcal L^{\mathrm{ext}}_\beta(n,\sigma)
:=
\inf_{I\in\mathcal I_{\beta,n,\sigma}}
\ \sup_{P\in\mathcal M_\beta(\sigma)}
\mathfrak r_{\mathrm{len}}(I,P).
\]
The constant interval procedure
\[
I_{\mathrm{full}}(z):=[-1,1]
\qquad(z\in\Omega_n)
\]
is measurable. Since \(|\theta(P)|\le1/2\) throughout the benchmark class,
\[
\{z:\theta(P)\in I_{\mathrm{full}}(z)\}=\Omega_n,
\qquad
\mathbb P_{P,n}\{\theta(P)\in I_{\mathrm{full}}\}=1.
\]
It is therefore honest, and
\[
\mathfrak r_{\mathrm{len}}(I_{\mathrm{full}},P)
=
\int_{\Omega_n}2\,d\mathbb P_{P,n}
=2.
\]
Consequently,
\[
0\le\mathcal L^{\mathrm{ext}}_\beta(n,\sigma)\le2<\infty,
\]
so this extended minimax value also admits an order-preserving conversion to a real number.
% lean: CausalSmith.Stat.RdTruesideNoiseFrontier.length_value_ne_top

\item
Fix any \(I\in\mathcal I_{\beta,n,\sigma}\), and define its two coverage events by
\[
\mathcal E_{\mathrm{cov},\pm}
:=
\{z\in\Omega_n:
\ell(z)\le\theta(P^\pm_{b,m})\le\upsilon(z)\}.
\]
These events are measurable. Since both alternatives are legal, honesty in \cref{obj:def:honest-class} supplies
\[
\mathsf P_+(\mathcal E_{\mathrm{cov},+})\ge\frac9{10},
\qquad
\mathsf P_-(\mathcal E_{\mathrm{cov},-})\ge\frac9{10}.
\]
The eventwise total variation bound gives
\[
\left|
\mathsf P_+(\mathcal E_{\mathrm{cov},+})
-
\mathsf P_-(\mathcal E_{\mathrm{cov},+})
\right|
\le
\operatorname{TV}(\mathsf P_+,\mathsf P_-)
\le\frac15,
\]
and therefore
\[
\mathsf P_-(\mathcal E_{\mathrm{cov},+})
\ge\frac9{10}-\frac15=\frac7{10}.
\]
Using the union--intersection identity under the probability law \(\mathsf P_-\), we obtain
\[
\begin{aligned}
\mathsf P_-\bigl(
\mathcal E_{\mathrm{cov},+}\cap\mathcal E_{\mathrm{cov},-}
\bigr)
&=
\mathsf P_-(\mathcal E_{\mathrm{cov},+})
+
\mathsf P_-(\mathcal E_{\mathrm{cov},-})
-
\mathsf P_-\bigl(
\mathcal E_{\mathrm{cov},+}\cup\mathcal E_{\mathrm{cov},-}
\bigr)
\\
&\ge
\frac7{10}+\frac9{10}-1
=
\frac35.
\end{aligned}
\]

On the intersection of the coverage events, both targets lie between \(\ell(z)\) and \(\upsilon(z)\), so
\[
\Delta_{\mathrm{wit}}
=
|\theta(P^+_{b,m})-\theta(P^-_{b,m})|
\le
\upsilon(z)-\ell(z)
\le
\max\{\upsilon(z)-\ell(z),0\}.
\]
Outside the intersection, its indicator is zero and interval length is nonnegative. Thus, on all of \(\Omega_n\),
\[
\Delta_{\mathrm{wit}}\,
\mathbf1_{\mathcal E_{\mathrm{cov},+}\cap\mathcal E_{\mathrm{cov},-}}(z)
\le
\max\{\upsilon(z)-\ell(z),0\}.
\]
Integration under \(\mathsf P_-\) now gives
\[
\begin{aligned}
\mathfrak r_{\mathrm{len}}(I,P^-_{b,m})
&\ge
\Delta_{\mathrm{wit}}\,
\mathsf P_-\bigl(
\mathcal E_{\mathrm{cov},+}\cap\mathcal E_{\mathrm{cov},-}
\bigr)
\\
&\ge
\frac35\,\Delta_{\mathrm{wit}}
=
\frac65\,\delta_{\mathrm{wit}}.
\end{aligned}
\]
% lean: CausalSmith.Stat.RdTruesideNoiseFrontier.twoPoint_interval_lower

\item
For every honest interval procedure, the risk at the minus alternative is bounded by its worst-case risk:
\[
\frac65\,\delta_{\mathrm{wit}}
\le
\mathfrak r_{\mathrm{len}}(I,P^-_{b,m})
\le
\sup_{P\in\mathcal M_\beta(\sigma)}
\mathfrak r_{\mathrm{len}}(I,P).
\]
Taking the infimum over honest procedures yields
\[
\mathcal L^{\mathrm{ext}}_\beta(n,\sigma)
\ge
\frac65\,\delta_{\mathrm{wit}}.
\]
Its finiteness then gives
\[
\mathcal L_\beta(n,\sigma)
\ge
\frac65\,\kappa\left(\frac{b}{m^2}\right)^\beta.
\]
Both arguments use the joint sample-and-seed laws, and hence apply to every allowed randomized procedure.
% lean: Causalean.Stat.le_minimaxValueENNReal
\end{enumerate}
\end{proof}
